\documentclass[12pt]{article}
\usepackage[brazil]{babel}
\usepackage[utf8]{inputenc}
\usepackage[T1]{fontenc}
\usepackage{amsmath, amssymb, amsthm}
\usepackage{geometry}
\usepackage{graphicx}
\usepackage{xcolor}
\usepackage{enumitem}
\usepackage{listings}

\definecolor{maincolor}{RGB}{0,100,50}
\usepackage{fancyhdr}
\definecolor{forestgreen}{RGB}{0,100,50}
\usepackage[most]{tcolorbox}

\lstset{
  basicstyle=\ttfamily\small,
  keywordstyle=\color{blue},
  stringstyle=\color{red},
  commentstyle=\color{green!50!black},
  breaklines=true,
  frame=single,
  numbers=left,
  numberstyle=\tiny\color{gray},
}
\newtcolorbox{exercisebox}[1]{
  title=\textbf{Exercício #1},
  fonttitle=\bfseries,
  coltitle=white,
  colbacktitle=maincolor
}
\newtcolorbox{block}[1]{%
  colback=green!5,
  colframe=maincolor,
  coltitle=white, 
  title=#1,
  fonttitle=\bfseries,
  boxrule=1.5pt,
  arc=3mm,
  left=6pt,
  right=6pt,
  top=6pt,
  bottom=6pt
}

\geometry{margin=2.5cm}
\pagestyle{fancy}
\fancyhf{}

\renewcommand{\headrulewidth}{2pt}
\renewcommand{\footrulewidth}{2pt}

\renewcommand{\headrule}{\hbox to\headwidth{%
  \color{forestgreen}\leaders\hrule height \headrulewidth\hfill}}

\renewcommand{\footrule}{\hbox to\headwidth{%
  \color{forestgreen}\leaders\hrule height \footrulewidth\hfill}}

\setlength{\headheight}{32pt}

\fancyhead[L]{%
  \includegraphics[height=1.2cm]{Universidade_de_Rio_Verde.jpg}
}

\fancyhead[R]{%
  \raggedleft
  \textcolor{forestgreen}{%
    \small
    \textbf{RESPOSTA - Lista de exercícios de Divisibilidade }\\
    Me.~Sergio Souza Novak\\
    \textit{Lógica para Engenharia de Software}
  }
}

\fancyfoot[C]{%
  \textcolor{forestgreen}{\thepage}
}

\begin{document}

% =====================================================
\begin{exercisebox}{1}
Prove que para quaisquer inteiros \(a,b,c\), se 
\(a \mid b\) e \(a \mid c\), então \(a \mid (b+c)\).
\end{exercisebox}

\begin{proof}
Sejam \(a,b,c \in \mathbb{Z}\). Suponha que \(a\) divide
\(b\) e que \(a\) divide \(c\). Pela definição de
divisibilidade, existem inteiros \(k_1\) e \(k_2\) tais que
\(b = ak_1\) e \(c = ak_2\).
Somando essas igualdades, obtemos
\(b+c = ak_1 + ak_2 = a(k_1+k_2)\).
Como o conjunto dos inteiros é fechado para a adição,
segue que \(k_1+k_2 \in \mathbb{Z}\).
Logo, existe um inteiro \(k\) tal que
\(b+c = ak\), concluindo-se que \(a\) divide \(b+c\).
\end{proof}

% =====================================================
\begin{exercisebox}{2}
Prove que, para quaisquer inteiros \(a\) e \(b\),
se \(a \mid b\), então \(a \mid (-b)\).
\end{exercisebox}

\begin{proof}
Sejam \(a,b \in \mathbb{Z}\) tais que \(a\) divide \(b\).
Então existe um inteiro \(k\) para o qual
\(b = ak\).
Multiplicando ambos os lados da igualdade por \(-1\),
obtemos \(-b = -(ak) = a(-k)\).
Como \(-k\) também é um inteiro,
segue diretamente da definição de divisibilidade que
\(a\) divide \(-b\).
\end{proof}

% =====================================================
\begin{exercisebox}{3}
Prove que, para quaisquer inteiros \(a,b,c\),
se \(a \mid b\) e \(b \mid c\),
então \(a \mid c\).
\end{exercisebox}

\begin{proof}
Sejam \(a,b,c \in \mathbb{Z}\).
Suponha que \(a\) divide \(b\) e que \(b\) divide \(c\).
Então existem inteiros \(k_1\) e \(k_2\) tais que
\(b = ak_1\) e \(c = bk_2\).
Substituindo a expressão de \(b\) na igualdade que define
\(c\), obtemos
\(c = (ak_1)k_2 = a(k_1k_2)\).
Como o produto de inteiros é ainda um inteiro,
temos \(k_1k_2 \in \mathbb{Z}\).
Portanto, pela definição de divisibilidade,
conclui-se que \(a\) divide \(c\).
\end{proof}

\end{document}